% Generated draft: requires mathematical and source review.
\documentclass[11pt]{article}
\usepackage[margin=25mm]{geometry}
\usepackage{fontspec}
\setmainfont{Latin Modern Roman}
\usepackage{amsmath,amssymb,mathtools,unicode-math}
\setmathfont{Latin Modern Math}
\usepackage{enumitem,needspace,xurl,hyperref}
\setlength{\parindent}{0pt}
\setlength{\parskip}{5pt}
\setlength{\emergencystretch}{3em}
\newcommand{\N}{\mathbb{N}}
\newcommand{\Z}{\mathbb{Z}}
\newcommand{\Q}{\mathbb{Q}}
\newcommand{\R}{\mathbb{R}}
\newcommand{\C}{\mathbb{C}}
\newcommand{\F}{\mathbb{F}}
\newcommand{\A}{\mathbb{A}}
\newcommand{\PP}{\mathbb{P}}
\newcommand{\E}{\mathbb{E}}
\newcommand{\eps}{\varepsilon}
\newcommand{\dd}{\,\mathrm{d}}
\begin{document}
\section*{UnsolvedMath ID: 800601}
\textbf{Source problem code:} catalogue entry
\subsection{Definitions and mathematical statement}
The supplied record asks whether
\[
\mathsf{U}_{\mathrm{fin}}(\mathcal O,\Omega)
=\binom{\Lambda}{\gimel(\Omega)}.
\]
Here equality is understood as equivalence of properties of spaces. The record does not specify the intended class of spaces or define the cover notation. The following gives an explicit, provisional countable open-cover interpretation; the ambient class and these conventions require confirmation against the cited source.

Let \(\mathscr K\) denote the intended class of spaces, which remains unspecified, and let \(X\in\mathscr K\). Under this interpretation, a cover of \(X\) is a countable family \(\mathcal U\) of open subsets of \(X\) such that \(\bigcup\mathcal U=X\) and \(X\notin\mathcal U\). Define:
\begin{itemize}
\item \(\mathcal O(X)\) is the family of these covers.
\item \(\Lambda(X)\) consists of those \(\mathcal U\in\mathcal O(X)\) for which every point of \(X\) belongs to infinitely many members of \(\mathcal U\).
\item \(\Omega(X)\) consists of those \(\mathcal U\in\mathcal O(X)\) for which every finite subset of \(X\) is contained in some member of \(\mathcal U\).
\item \(\Gamma(X)\) consists of the infinite \(\mathcal U\in\mathcal O(X)\) for which every point of \(X\) belongs to all but finitely many members of \(\mathcal U\).
\end{itemize}
For a family \(\mathscr B(X)\) of covers, the property \(\mathsf U_{\mathrm{fin}}(\mathcal O,\mathscr B)\) means that for every sequence \((\mathcal U_n)_{n\geq1}\) in \(\mathcal O(X)\), each \(\mathcal U_n\) having no finite subcover of \(X\), there are finite subfamilies \(\mathcal F_n\subseteq\mathcal U_n\) such that
\[
\left\{\bigcup\mathcal F_n:n\geq1\right\}\in\mathscr B(X).
\]
A cover \(\mathcal U\) is called multifinite if it has a partition \((\mathcal P_n)_{n\geq1}\) into nonempty finite subfamilies satisfying \(\bigcup\mathcal P_n=X\) for every \(n\). In the provisional interpretation, \(\gimel(\mathscr B)(X)\) consists of covers \(\mathcal U\) that are multifinite or admit a partition \((\mathcal P_n)_{n\geq1}\) into nonempty finite subfamilies for which
\[
\left\{\bigcup\mathcal P_n:n\geq1\right\}\setminus\{X\}
\in\mathscr B(X).
\]
The notation \(\binom{\Lambda}{\gimel(\mathscr B)}\) means that every \(\mathcal U\in\Lambda(X)\) has a subfamily \(\mathcal V\subseteq\mathcal U\) belonging to \(\gimel(\mathscr B)(X)\).

With these conventions, the question is whether
\[
\forall X\in\mathscr K,\qquad
X\models\mathsf U_{\mathrm{fin}}(\mathcal O,\Omega)
\quad\Longleftrightarrow\quad
\forall\mathcal U\in\Lambda(X)\ \exists\mathcal V\subseteq\mathcal U:
\mathcal V\in\gimel(\Omega)(X).
\]
This is a conditional formalization of the displayed notation: the supplied extract does not determine \(\mathscr K\), and no assertion over all topological spaces is substituted for that missing scope.

As background, the extract states
\[
\mathsf U_{\mathrm{fin}}(\mathcal O,\Gamma)
=\binom{\Lambda}{\gimel(\Gamma)}
\]
with reference label \texttt{hureslaloms}. It also attributes to Zdomskyy a proof that NCF implies a positive answer to the question involving \(\Omega\). These are supplied background assertions, not independently verified results here. The extract does not define NCF or provide the scope and hypotheses of that implication; NCF is not imposed as a hypothesis of the question itself.
\subsection{Short English statement}
The question asks whether, for every space in the intended class, \(\mathsf U_{\mathrm{fin}}(\mathcal O,\Omega)\) is equivalent to every large open cover having a subcover in \(\gimel(\Omega)\). The definitions above give a provisional countable open-cover reading. The intended class of spaces and the original conventions remain unresolved. The registry labels the problem open; this draft does not independently verify that status.
\subsection{Sources}
\begin{itemize}
\item Local UnsolvedMath registry, supplied complete record with permanent ID \(800601\), source problem code \texttt{catalogue entry}, title ``Selection Principles in Mathematics --- Problem 29.'' The registry attributes the proposal to Tsaban, with proposed year \(2006\), and supplies status \texttt{open}. Its update timestamp is \texttt{2026-10-16T14:00:00Z}; this is registry metadata, not independent verification of mathematical status.
\item Tsaban, \emph{Selection Principles and Special Sets of Reals: Open Problems} (2006), Problem 29. Supplied source URL: \url{https://arxiv.org/abs/math/0609601}. The provenance reports access on \texttt{2026-10-16} and extraction method \texttt{arxiv\_latex}. That access date is supplied provenance; no source visit or independent verification is claimed here.
\item The extracted statement cites \texttt{hureslaloms} for the background equality involving \(\Gamma\). The supplied record does not resolve this reference label to a bibliographic entry.
\item The extracted statement attributes to Zdomskyy the implication from NCF to a positive answer. No separate bibliographic citation for that result is supplied.
\end{itemize}
\end{document}
